\chapter{Standards, industries and financing instruments}
\label{ch:modules}

\section{A common core with distinct accounting books}

The common core governs evidence, time, money, events, journals and provenance.
Reporting standards govern recognition, measurement and presentation within
that core. The design preserves three levels of meaning: economic terms such
as a promise to pay cash; an accounting treatment under a particular standard;
and a company's reported label or extension. These levels are connected by
explicit rules and mappings, rather than collapsed into one universal chart
of accounts.

A rule package includes a standards register with jurisdiction, issuer class,
effective dates, transition provisions, permitted elections, paragraph
references, implementation version and tests. The package selector uses the
entity, reporting basis, period and adopted elections. It must distinguish a
standard's mandatory effective date from permitted early adoption. For example,
the IFRS Foundation's IFRS 18 overview states an effective date of 1 January
2027 with earlier application permitted \citep{ifrs18}. A 2026 historical
report cannot be assigned that presentation regime solely because the latest
installed package includes it.

IFRS, US GAAP and ASBE packages share economic contracts and events where
appropriate, but produce separate accounting books. A reconciliation between
books consists of explained adjustments with origins in differing treatments,
not a renaming of accounts. The first implementation can preserve all three
source bases even while only a selected subset has executable rules. Such a
source is represented as observed information with unsupported simulation
treatments, rather than being silently converted to the supported basis.

The official IFRS overviews for IFRS 9, IAS 32 and IFRS 17 establish the relevant
scope of financial-instrument measurement, presentation and insurance
contracts \citep{ifrs9,ias32,ifrs17}. They are orientation sources, not substitutes
for reviewed paragraph-level implementations. Public retrieval of the selected
Chinese Ministry of Finance page and SEC API documentation failed during this
campaign; the source register records these access gaps. ASBE and US GAAP
semantic modules therefore have explicit primary-text acquisition and
accounting-review tasks before they can be marked executable and compliant.

\section{Non-financial companies}

The initial industry package covers sales, collections, purchases, payments,
inventory movements, payroll, capital expenditure, depreciation and basic
financing. Revenue recognition is a separate rule from invoicing and cash
collection. A service contract may generate a performance event before an
invoice, creating a different asset from an unconditional receivable. Customer
advances create obligations before revenue. Returns, rebates and variable
consideration need separate assumptions and subsequent adjustments.

Inventory state includes quantity, cost, location, ownership and lot or cohort.
A costing method is selected within the reporting package's permissible
domain. An aggregate stock can support a simple model only if its approximation
is declared. A commodity-price shock can expose differences between old and
new inventory that an average value suppresses. Fixed-asset state likewise
requires acquisition dates, components, useful lives, residual values,
depreciation method, tax basis and impairment history when those affect the
requested forecast.

Industry extensions add economic mechanisms without bypassing the accounting
core. Construction requires performance and completion states; extractive
industries require reserves, depletion and restoration obligations; utilities
may require rate-regulation treatment; software businesses may require
contract liabilities and development-cost decisions. The module manifest
states what is represented, what is executable, and what remains aggregated.

\section{Debt and hybrid financing}

The contract layer for a debt instrument stores issuer and holder perspective,
currency, issue and settlement dates, issue proceeds and legal face amount.
Its payment terms include coupon/index/spread, reset and fixing dates,
day-count convention, payment calendar, amortisation schedule and maturity.
Its rights and protections include seniority, collateral, covenants, options,
conversion terms, guarantees and linked hedges. Missing terms are
explicit unknowns. A single carrying amount cannot stand in for this record.

The accounting layer stores the classification decision and its evidence,
initial measurement, carrying amount, accrued interest, effective-rate method,
unamortised fees or discounts, fair-value disclosures, impairment or credit
adjustments where applicable, and presentation mappings. Contractual interest,
effective-interest expense, realised cash payments and fair-value changes are
separate quantities. The zero-coupon example in Chapter~\ref{ch:execution}
provides a minimal conformance case for this separation.

The lifecycle must support issuance, recognition, accrual, reset, coupon,
principal repayment, modification, derecognition, default, recovery and
conversion. A modification event first determines which accounting operation
the selected standard requires; it does not always mean issuing a new loan.
An embedded option can affect valuation and classification without being
exercised. A derivative linked as a hedge requires both the economic contract
and any applicable accounting designation; economic risk reduction alone does
not establish hedge-accounting eligibility.

Hybrid instruments make interpretation particularly consequential. IAS 32's
overview discusses liability/equity classification and the separation of a
compound instrument \citep{ias32}. The executable classification must consider
the full terms and the selected standard version. US GAAP or ASBE treatment
must be implemented under its own rules rather than inferred from an IFRS
classification. Where the terms are not disclosed sufficiently, alternative
classifications can coexist only if each is actually compatible with the
known terms and applicable rules.

\section{Equity and ownership}

Equity has both monetary accounts and security quantities. The security
register records share class, issued and outstanding quantities, treasury
holdings, par or stated value where applicable, voting rights, dividend rights,
redemption features, options and conversion relationships. Monetary equity
accounts include contributed capital, reserves, retained earnings and other
components required by the selected basis. These records evolve together
through issuance, repurchase, cancellation, split, dividend and conversion
events.

A stock split changes quantities and per-share measures without necessarily
changing aggregate monetary equity. A cash dividend affects distributions and
cash rather than current operating expense under the corresponding ordinary
equity treatment. Repurchases and employee share awards require their own
rules. Earnings per share uses a dated quantity schedule and rules for potential
ordinary shares; simply dividing profit by the closing share count is
insufficient when shares change during the period.

Funding simulations must connect the quantity register to market assumptions.
An equity issue at a scenario-dependent price changes cash raised, shares and
dilution. The model must state whether it targets a fixed cash amount, a fixed
share count or a capital ratio. These are different management decisions even
when the historical statement would report the same aggregate proceeds.

\section{Banks: accounting capital is not regulatory capital}

A bank package needs instrument and cohort states for loans, securities,
deposits, wholesale funding, repos, derivatives, collateral and commitments.
For interest-rate analysis, currency and repricing buckets matter alongside
remaining maturity. For credit analysis, exposure, default state, loss severity,
recoveries, guarantees and their dependencies matter. For liquidity analysis,
deposit behaviour, collateral eligibility, encumbrance and payment timing matter.
These states can be represented at an explicitly chosen aggregation level;
the public record rarely identifies every contract.

Loan origination can debit loans and credit a customer deposit in the bank's
book. Subsequent transfer to another bank produces settlement flows. Treating
all loan growth as an immediate cash outflow would therefore encode a specific
funding assumption, not a universal banking identity. A system-wide model must
also link counterparties and settlement assets; single-company simulation
records external counterparties without claiming to close the whole economy.

An allowance calculation is a dated measurement based on forecast information.
For a single horizon, let EAD be a fixed exposure at default, PD the probability
of default, and LGD the mean loss fraction conditional on default. Ignoring
discounting, the illustrative expected loss in currency units is
\[
 L=EAD\times PD\times LGD.
\]
With exposure 100, default probability 0.02 and loss given default 0.4, it
produces 0.8. Changing those inputs to 0.04 and 0.5 produces 2, so the allowance
increase is 1.2 if there are no other movements. This calculation is a
demonstration model. A standard-specific expected-credit-loss implementation
must determine horizon, staging or other classification, discounting,
forward-looking information, write-offs and purchased-credit-impaired or
other special cases. The simple formula is not an IFRS 9 or US GAAP compliance
implementation.

Regulatory capital is computed in a separate prudential layer with its own
perimeter, eligibility, deductions, risk weights and transition provisions.
Accounting equity feeds a reconciliation to capital; it is not automatically
Common Equity Tier 1. Liquidity and capital constraints influence the
management policy, which in turn generates asset sales, funding or capital
events. A breached ratio can be an economically meaningful scenario outcome.
No hidden capital injection is permitted to repair it.

\section{Insurance: cohorts and future service}

An insurer requires policy cohorts or groups with issue dates, coverage,
premiums, benefits, expenses, lapses, claims development, guarantees and
reinsurance relationships. Asset portfolios need their own contracts and
valuations. The liability state must distinguish expected future cash flows,
discounting, risk adjustments or margins, unearned profit measures where
applicable, incurred claims and movements between opening and closing states.
The IFRS 17 overview's distinction between measuring future cash flows and
recognising profit as service is provided motivates this separation
\citep{ifrs17}.

The module API accepts policy/cohort state, economic assumptions, non-economic
assumptions and information date. It returns a measurement vector, movement
attribution and proposed recognition events. Experience differences, changes
in assumptions, finance effects and service effects remain separately
identifiable. Reinsurance held has its own record and measurement; it is not
an automatic percentage offset to every insurance liability.

Scenarios can alter discount curves, mortality, morbidity, lapse rates,
inflation, catastrophe losses and policyholder behaviour. Their joint
dependence and cohort exposure must be declared. A portfolio total can match
history while hiding very different guarantee and duration profiles. This
makes alternative economic histories especially valuable, but also limits
what public-data reconstruction can identify. A full insurance package needs
actuarial and accounting review beyond the core release.

\clearpage
\section{Capability contracts for sector packages}

\begin{longtable}{L{26mm}L{55mm}L{55mm}}
\toprule Package & Required state and operations & Minimum acceptance evidence\\\midrule\endhead
Core corporate & Working-capital bridges, lots/assets, ordinary debt and equity,
period closing and cash-flow mappings & Exact synthetic cases, reviewed public
filing reconstruction, all statement reconciliations.\\
Financial instruments & Terms, calendars, resets, options, measurement
categories, modifications and hedges & Contract schedule comparisons, valuation
checks, paragraph-based posting cases and transition cases.\\
Banking & Credit cohorts, deposits, collateral, liquidity, regulatory capital,
management actions & Allowance roll-forwards, repricing and runoff stresses,
capital reconciliation and default/funding-failure cases.\\
Insurance & Policy groups, future cash flows, service and finance components,
reinsurance and asset interaction & Cohort movement reconciliation, assumption
change cases, service recognition and reviewed actuarial calculations.\\
Multi-basis reporting & Versioned IFRS/US GAAP/ASBE rules, entity elections and
transition treatment & Matched economic-event cases with explained basis
differences and independent accounting review.\\\bottomrule
\end{longtable}

Every package declares capabilities at operation level. ``Supports bonds'' is
too vague; ``supports fixed-rate, single-currency, non-callable amortised-cost
issuance and redemption under this reviewed rule set'' is implementable and
testable. Unimplemented operations retain their source and contract
representation while producing an explicit capability error when execution
requires them.

\section{A matched inventory case across reporting profiles}

Standards packages need operation-level decisions. For interchangeable goods,
the inspected AASB~102 text incorporating IAS~2 permits FIFO or weighted
average in paragraph 25 and explains both in paragraph 27
\citep{aasb102}. The reproduced ASBE~1 text specifies FIFO, weighted average
or specific identification in Article 14; its official Ministry version could
not be retrieved in this review, so that reproduction has provisional source
status \citep{asbe1copy}. Caterpillar's company-authored 2024 report states
that it principally uses LIFO and supplies a FIFO comparison
\citep[Note 1.C]{caterpillar2024}. That is evidence of a disclosed US-GAAP
company policy; full ASC~330 paragraph authority remains an acquisition task.
It is not a rule that every US company must use LIFO.

Take the same purchases in three separate books: ten units at 5, followed by
ten at 7, then delivery of twelve units. FIFO charges $10(5)+2(7)=64$ to
cost of sales and leaves inventory 56. LIFO charges $10(7)+2(5)=80$ and
leaves inventory 40. Periodic weighted average uses cost $120/20=6$, charges
72 and leaves 48. The physical deliveries and supplier cash flows agree;
profit and carrying inventory differ by the selected costing rule. A concept
mapping cannot remove that difference. The FIFO/average selector in an
IAS~2 profile rejects LIFO as an unsupported election under the inspected
paragraph. The US company profile admits its evidenced LIFO policy while
retaining its narrower authority. An ASBE package remains source-provisional
until the attributed text and applicable amendments are officially verified.

Inventory impairment illustrates another required memory. Under inspected
AASB~102 paragraphs 28--34, an item costing 56 written down to net realisable
value 50 records expense 6. If the conditions reverse and net realisable
value becomes 54, reversal is 4, capped by original cost; the previous
write-down must therefore remain in state. The reproduced ASBE~1 Article 19
also addresses reversal within the earlier provision. This paragraph is not
used to assert the US treatment, whose primary authority was unavailable.
Source limitations constrain the treatment claim, while the separate-book
method and exact costing calculations remain fully specified.

\section{A compound bond exposes classification dependencies}

Consider an issuer receiving 100 for a two-year bond paying annual coupons of
5 and principal 100 at maturity, with a holder option to convert that principal
into ten ordinary shares. Assume its currency is the issuer's functional
currency, no other options, fees or contingent settlements, and that the
conversion feature qualifies as equity under the selected rule. A comparable
liability without conversion yields 10\%. The inspected IAS~32-equivalent
text requires separate liability and equity components and measures equity
as the residual \citep[paragraphs 28--32, AG30--AG32]{aasb132}.

The liability's present value is
\[
 L_0=\frac{5}{1.1}+\frac{105}{1.1^2}
     =\frac{11050}{121}\approx91.3223,
 \qquad E_0=100-L_0\approx8.6777.
\]
Issue debits cash 100 and credits these two components. Under the stated
effective-interest schedule, the first coupon leaves liability
$1.1L_0-5=1050/11\approx95.4545$; after the second coupon it is 100.
Redemption debits liability and credits cash 100. Alternatively, conversion
at maturity debits liability and credits equity 100 while issuing ten shares;
the original equity component remains in equity. AG32 supports the absence
of a gain or loss on that maturity conversion. Exact rational valuations are
kept before the declared monetary rounding.

A variable-number share settlement, cash alternative or other clause can
change classification. The contract record must therefore retain these terms;
it cannot select liability/equity accounts independently of the coupon,
conversion and settlement rules. Early repurchase and induced conversion
invoke separate operations: the inspected AG33--AG35 text describes allocation
and the induced-conversion consideration difference. They are not implemented
by changing the maturity date of the simple example. US-GAAP and ASBE
classification require their own dated authority; the IFRS-oriented split is
never silently reused under another basis.

\section{Control and an elimination book}

The consolidation procedure selects the control graph at each date, aligns
supported accounting policies and presentation currencies, combines controlled
entity books, and posts acquisition, elimination and non-controlling-interest
adjustments in a separate book. A control conclusion needs evidence and a
rule version. Ownership percentage alone is insufficient. Unknown purchase
allocation, currency translation or unmatched intragroup differences remain
named unresolved adjustments rather than unexplained balancing entries.

For a simplified acquisition, assume control is established, all net assets
are already at the relevant acquisition measurements, no tax or goodwill
arises, and non-controlling interest is measured proportionately. A parent
has cash 60, investment 40 and equity 100; an 80\%-owned subsidiary has cash
30, plant 20 and equity 50. Consolidation debits subsidiary equity 50,
credits investment 40 and credits non-controlling interest 10. Group assets
are cash 90 plus plant 20, and group equity is 100 attributable to the parent
plus 10 to other owners.

A subsequent intragroup loan of 10 shifts cash from parent to subsidiary,
creates a receivable and payable, and is eliminated by debiting payable and
crediting receivable 10. Group totals remain 110. An unmatched loan currency
or settlement date would block that exact elimination. If the parent instead
sells its entire interest for 44 immediately after the initial state and
loses control, the group removes subsidiary cash 30, plant 20 and
non-controlling interest 10, recognises cash consideration 44 and gain 4;
remaining cash/equity is 104. More complex acquisitions and retained interests
require additional measurement operations; this example fixes the structure
and arithmetic without certifying a complete consolidation standard.

\section{A bank's cash, allowance and capital in one path}

The bank reference state contains settlement cash 30, gross loans 100,
allowance 0.8, deposits 110 and equity 19.2. Originating a loan of 20 first
debits loans and credits the borrower's deposit. A subsequent transfer of 15
to another bank debits deposits and credits settlement cash. Loan origination
and external cash use are consequently distinct events. Receive loan interest
6 in cash and credit deposit interest 2 to customer balances. Remeasure the
allowance to 2.4, posting an additional loss of 1.6. The resulting state is
cash 21, gross loans 120, allowance 2.4, deposits 117 and equity 21.6.
Net assets are $21+120-2.4=138.6=117+21.6$.

Now 25 of deposits runs off. The declared funding policy uses cash first,
then a committed facility with limit 10. It draws 4 and settles 25, leaving
cash zero, deposits 92 and wholesale debt 4. A runoff of 40 instead leaves
9 unpaid even after the maximum facility draw. The simulator records default
and retains the unpaid deposit obligation. It does not assume unlimited
central-bank or shareholder funding.

The prudential calculation is a separate, illustrative policy. Deduct 2 from
accounting equity to obtain capital 19.6; apply a 75\% risk weight to gross
loans for risk-weighted assets 90. The resulting ratio is $19.6/90$.
Neither the deduction nor the weight is a Basel default. Actual prudential
packages require jurisdiction, exposure class, eligibility, buffers and dated
rules, with a reconciliation from the accounting book. The structure follows
the need to model capital and liquidity jointly in
\citet[Section 2]{behn2019}, without transferring that paper's estimated
preferences or shock parameters to this bank.

For a richer credit adapter, partition borrower cohorts by currency, vintage,
credit state, remaining maturity, collateral and measurement category. At date
$t$, supply scenario probabilities, conditional default-time probabilities,
exposure paths, conditional loss severity and discount factors. Its expected
loss calculation has the declared form
\[
 L_t=\sum_s p_s\sum_{k\in H}
  d_{t,k}\Pr(\tau=k\mid s,\mathcal F_t)
  \E[\mathrm{EAD}_k\mathrm{LGD}_k\mid\tau=k,s,\mathcal F_t].
\]
The conditional product expectation retains dependence between exposure and
severity. The horizon $H$ concerns eligible default times; associated cash
shortfalls may continue later. The rule package decides twelve-month versus
lifetime default scope, credit-state transfers, interest basis and write-off
recognition. Changes in allowance must be attributed to origination,
derecognition, transfers, revised expectations, currency and write-offs.
Contract and prudential support can be specified now, while full IFRS~9/CECL/
ASBE classification and paragraph-level validation remain later obligations.

\section{Insurance service and reinsurance in a complete cohort}

The example uses one-year policies and a one-year proportionate treaty starting
when those policies are recognised, with uniform service for each. Both groups
qualify for premium allocation through the one-year conditions in paragraphs
53(b) and 69(b); paragraphs 62--62A determine treaty recognition. Assume no
acquisition cash flows, investment components, significant financing or onerous
remaining coverage. Claims and recoveries are expected within a year of the
claims being incurred. The inspected AASB~17 text incorporating IFRS~17 specifies
the movements in paragraphs 53--59, allocation in B126 and separate reinsurance
requirements in 63--70A \citep{aasb17}. Gross and ceded groups retain distinct
identities.

Begin with capital and cash of 20. Receive premium 120 and recognise a
remaining-coverage liability of 120. Pay reinsurance premium 24 and recognise
a coverage asset of 24. After half the uniform service, recognise revenue 60
by releasing gross coverage and reinsurance expense 12 by releasing its
coverage asset. Estimated incurred claims are 40, of which 30 is paid;
the remaining-claim risk adjustment is assumed to be 2. Gross claims expense
is therefore 42 and incurred-claims liability is 12. Assume the treaty
reimburses half the gross expected claims, recovery cash of 15 has arrived,
the risk transferred is valued at 1, and reinsurer non-performance has zero
expected effect in this example. The reinsurance claims asset is
$20+1-15=6$, and recovery income is 21.

Cash is $20+120-24-30+15=101$. Assets are cash 101, remaining reinsurance
coverage 12 and claims recovery 6, totalling 119. Liabilities are remaining
gross coverage 60 plus incurred claims 12, totalling 72. Profit is
$60-42-12+21=27$ and equity is $20+27=47$. These balances reconcile without
netting reinsurance against the insurance liability. If counterparty credit
risk changes, the recovery measurement changes through its own inputs; it
does not mechanically halve every gross liability or risk adjustment.

Longer-duration business adds a fulfilment-cash-flow vector, discount curve,
risk adjustment, contractual service margin, loss component and coverage-unit
schedule. Select portfolio and annual cohort, assess onerous grouping, project
cash flows within contract boundaries, discount with the declared financial
assumptions, measure non-financial risk separately, and attribute subsequent
movements to new business, experience, assumptions, finance and service.
Future-service changes adjust the supported service-margin operation subject
to the loss rules; current/past-service effects enter their specified result
components. Unsupported variable-fee, transition or loss-recovery operations
must stop those dependencies explicitly. A PAA example does not implement
these longer-duration methods by implication.

\enlargethispage{3\baselineskip}
An actuarial reserve adapter also needs a concrete estimation contract.
For cumulative claims $C_{i,k}$ by origin year $i$ and development year $k$,
the chain-ladder point method estimates each factor from matched observed
cells, $\widehat f_k=\sum_i C_{i,k+1}/\sum_i C_{i,k}$, then projects the
latest cumulative claim along the remaining factors. A triangular example
with rows $(100,150,180)$, $(120,180)$ and $(140)$ gives factors 1.5 and
1.2, outstanding amounts 36 and 112 for the unfinished rows, and total 148.
\citet[Sections 2--3, equations (1)--(3)]{mack1993} makes explicit the
conditional-mean and variance assumptions and independence across origin
years; its aggregate-error discussion retains dependence from shared estimated
factors. A single final-development observation cannot supply an empirical
tail variance without an additional assumption. The adapter therefore returns
the point projection and marks uncertainty estimation unavailable in this tiny
example. It supplies neither a full predictive distribution nor an IFRS risk
adjustment. Calendar effects, inflation, changing claims handling and tail
selection require separate sensitivity and validation before real use.
